% Part: sets-functions-relations
% Chapter: functions
% Section: isomorphisms
%READERNOTE{SOL6-A340: समरूपतेच्या व्याख्येतील द्विदिशा स्पष्ट केली; सर्व संबंध आणि उदाहरणे जपली.}

\documentclass[../../../include/open-logic-section]{subfiles}

\begin{document}

\olfileid{sfr}{fun}{iso}
\olsection{समरूपता}

\begin{explain}
\emph{समरूपता} म्हणजे संबंधित संचांची रचना जपणारे एकास-एक
आच्छादन. ही रचना संचांतील !!{element}s मधील संबंधांवर अवलंबून
असते. पुढील दोन संच पाहा: $X=\{1,2,3\}$ आणि $Y=\{4,5,6\}$.
दोन्ही संचांवर उत्तरवर्ती, लहान असणे आणि मोठे असणे हे संबंध
आहेत. या दोन संचांमधील समरूपता म्हणजे या रचना जपणारे
!!a{bijection}. म्हणून $f \colon X
\to Y$ हे !!a{bijective} फलन समरूपता असेल, जर
$i<j$ तेव्हा आणि केवळ तेव्हा $f(i)<f(j)$, $i>j$ तेव्हा आणि केवळ तेव्हा $f(i)>f(j)$, आणि $j$ हा $i$ चा उत्तरवर्ती असेल
तेव्हा आणि केवळ तेव्हा $f(j)$ हा $f(i)$ चा उत्तरवर्ती असेल.
\end{explain}

\begin{defn}[समरूपता]
$U$ ही जोडी $\langle X, R\rangle$ आणि $V$ ही जोडी $\langle
Y, S\rangle$ असो; येथे $X$ आणि $Y$ हे संच, तर $R$ आणि $S$ हे
अनुक्रमे $X$ आणि $Y$ वरील संबंध आहेत. $X$ पासून $Y$ कडे जाणारे
!!{bijection}~$f$ हे $U$ पासून $V$ कडे जाणारी \emph{समरूपता}
असेल तेव्हा आणि केवळ तेव्हा ते संबंधात्मक रचना जपते; म्हणजे
$X$ मधील कोणत्याही $x_{1}$ आणि $x_{2}$ साठी
$\tuple{x_1,x_2} \in R$ तेव्हा आणि केवळ तेव्हा
$\tuple{f(x_1),f(x_2)} \in S$.
\end{defn}

\begin{ex}
$X=\{1,2,3\}$ आणि $Y=\{4,5,6\}$ हे दोन संच, तसेच
लहान असणे आणि मोठे असणे हे संबंध पाहा. $f(x) = 7-x$ असे
असलेले $f\colon X \to
Y$ हे फलन $\tuple{X,<}$ आणि $\tuple{Y,>}$ यांमधील
समरूपता आहे.
\end{ex}

\end{document}
